(1) 切向波数 \(k_x=k_0n_1\sin\theta\) 在各平面界面连续。低折射率间隙纵向波数满足 \(k_{z2}^2=k_0^2n_2^2-k_x^2=-\kappa^2\)，因此场由 \(e^{+\kappa z}\)、\(e^{-\kappa z}\) 的线性组合构成，不能将它们当作在无限间隙中携带独立远场传播功率的普通行波。左右高折射率区域仍为真实传播波。

(2) TE 电场 \(E_y\) 及其法向导数（非磁性介质）在每一界面连续。将间隙传输矩阵 \(\begin{pmatrix}\cosh\kappa d&\sinh\kappa d/\kappa\\\kappa\sinh\kappa d&\cosh\kappa d\end{pmatrix}\) 与两侧传播边界连接，得到透射幅值的模平方 \(T=[1+((k_z^2+\kappa^2)^2/(4k_z^2\kappa^2))\sinh^2(\kappa d)]^{-1}\)。两侧介质相同，所以功率透射率等于透射电场振幅模平方，无额外不同光学导纳比因子。

(3) 当 \(d\to0\)，\(\sinh\kappa d\to0\)，所以 \(T\to1\)，对应两块高折射率介质重新接触。厚间隙 \(\kappa d\gg1\) 时 \(\sinh^2\kappa d\sim e^{2\kappa d}/4\)，于是 \(T\sim[16k_z^2\kappa^2/(k_z^2+\kappa^2)^2]e^{-2\kappa d}\)。指数衰减以强度计的尺度为 \((2\kappa)^{-1}\)，场幅穿透深度则为 \(\kappa^{-1}\)，必须区分两者。

(4) 无吸收、无外部散射通道的静态双界面结构由能流连续性给 \(R+T=1\)。单界面全反射虽无净透射到无限低折射率区，但有限间隙的两个倏逝解经第二界面可耦合出右侧传播波，故 \(T\) 可非零。极限 \(\kappa\to0\) 与厚间隙近似不交换，靠近临界角必须使用完整双曲函数公式，而不能套 \(e^{-2\kappa d}\) 厚势垒公式。
